% Lösungen Einheit 5 von 5 – Das Integral als Flächenbilanz (zusammenbau v0.1)
\einheitenkopf{Einheit 5 von 5 · Das Integral als Flächenbilanz}

\erg{17}{a) Wert negativ – nicht die Fläche, die Fläche ist sein Betrag \quad b) Wert $= 8$ (Trapez unter dem Graphen, ganz oberhalb der Achse) \quad c) Das Stück links liegt über, das Stück rechts unter der x-Achse; beide sind wegen der Punktsymmetrie zum Ursprung gleich groß, ihre orientierten Inhalte heben sich auf. \quad d) $12$: der ungerade Anteil liefert auf dem zu null symmetrischen Intervall null, die Konstante ein Rechteck $3 \cdot 4$. \quad e) kleiner – das Stück von $3$ bis $5$ zählt negativ \quad f) Trapez $= 6$; genau $\frac{14}{3} \approx 4{,}67$; Überschätzung, weil der Graph unter der Sehne liegt \quad g) Mittelwert $= \frac{8}{2} = 4$ – das Rechteck der Höhe $4$ über dem Intervall hat denselben Inhalt wie die Fläche unter dem Graphen \quad h) Richtig: Beide Faktoren sind nie negativ, das Quadrat nicht und die e-Funktion nicht; der Graph liegt also nie unter der x-Achse, und das Integral zählt kein Stück negativ.}
\erg{18}{a) Fehler: Integralwert und Fläche gleichgesetzt, über die Nullstelle $0$ hinweg. Richtig: $A = 4 + \frac{1}{4} = 4{,}25$ \quad b) Der Graph ist punktsymmetrisch zum Ursprung: Jedem Flächenstück rechts entspricht links ein gleich großes auf der anderen Seite der x-Achse, die orientierten Inhalte heben sich auf. \quad c) Integral $= 72$ kWh; Mittelwert $= 72 : 12 = 6$ kW – eine Anlage mit konstant $6$ kW liefert in $12$ Stunden dieselbe Energie \quad d) Parallele bei $x = 2$ (Trapez $\frac{1 + 2}{2} \cdot 2 = 3$); Term: Integral von $0$ bis $2$ über $f(x)$}

18d)

\begin{ksys}[xmin=-1,xmax=5,ymin=-1,ymax=5,klein]\flaeche{0.5*\x+1}{0}{2}{A} \funktion{0.5*\x+1}{f}\end{ksys}

